% Timing diagram of a result: following one result of job module 12 from the
% job module to the write to the display memory. Rendered to
% timing_result.png by "make diagrams" (see the DIAGRAMS variable in the
% top-level Makefile). It is shown in ALGORITHM.md, in the section "Timing of
% a job and a result" of the Dispatcher, next to the table of the clock cycles
% of a result, and the clock cycles r, c, ... are the same as in the table.
%
% All values are read off a simulation of the first instance of the
% dispatcher in sim/dispatcher_tb.vhd (i_dispatcher: 16 job modules in groups
% of 5, so job module 12 is in group 2 and res_scheduler has five positions;
% 64 picture columns with an address distance of 512; one pixel in each
% write), in the first picture: r is the clock cycle from 645 ns to 655 ns (the
% rising edge of clk at 645 ns). The result is the first row (0) of the job of
% picture column 6, with the count 18, so it is written to address 6*512+0 =
% 3072. Here job module 12 is selected as early as possible, c = r+4. Anything
% that changes the processes or the registers of the dispatcher, or that
% testbench, has to change this file too.
%
% Drawing convention, the same as in timing_job.tex: single bits are always
% drawn with their true value. A bus is white in the clock cycles where it
% carries the result of job module 12, light grey with its value where it
% carries another result, and grey without a value where it carries nothing
% (its valid flag is low).
%
% Many job modules have a result at the same time, so res_scheduler selects
% other job modules in most clock cycles, and the steps of their results
% overlap with those of this one. The counter of res_scheduler is at group 2
% (job module 12) in r+3, and job module 12 is the candidate of group 2 in that
% clock cycle (cand_r, not drawn), so it is selected in c = r+4. The job module
% starts the next row (res_start_r) as soon as it sees the acknowledge in
% c+2.

\documentclass[border=8pt]{standalone}

\usepackage[T1]{fontenc}
\usepackage{helvet}
\renewcommand{\familydefault}{\sfdefault}

\usepackage{tikz-timing}

% White: the result of job module 12. Light grey: another result. Grey (U):
% nothing.
\tikzset{
   timing/d/background/.style={fill=white},
   timing/u/background/.style={fill=black!25},
   wh/.style={timing/d/background/.style={fill=white}},
   gr/.style={timing/d/background/.style={fill=black!8}},
   timing/name/.style={font=\ttfamily\small},
   timing/d/text/.style={font=\scriptsize},
}

\begin{document}
\begin{tikztimingtable}[timing/slope=0.08, timing/coldist=6pt, xscale=2.2, yscale=1.4]
clk                                  & 11{HL} \\
res\_valid\_i(12)                    & 2L 14H 6L \\
res\_addr\_i(12)                     & 2U [wh] 14D{0} 6U \\
res\_data\_i(12)                     & 2U [wh] 14D{18} 6U \\
res\_busy\_r(12)                     & 4H 12L 6H \\
res\_ready\_s(12)                    & 4L 10H 8L \\
res\_scheduler/cnt\_r                & [gr] 2D{3} 2D{4} 2D{0} 2D{1} [wh] 2D{2} [gr] 2D{3} 2D{4} 2D{0} 2D{1} 2D{2} 2D{3} \\
idx\_valid\_r                       & 4H 2L 6H 4L 6H \\
idx\_iterator\_r                     & [gr] 2D{10} 2D{15} 2U 2D{0} 2D{9} [wh] 2D{12} 4U [gr] 2D{1} 2D{5} 2D{13} \\
grp\_valid\_r                        & 6H 2L 6H 4L 4H \\
acc\_job\_addr\_r                    & [gr] 2D{2} 2D{4} 2D{9} 2U 2D{10} 2D{3} [wh] 2D{6} 4U [gr] 2D{11} 2D{15} \\
res\_ack\_o(12)                      & 14L 2H 6L \\
acc\_col\_dd                         & 4U [gr] 2D{1024} 2D{2048} 2D{4608} 2U 2D{5120} 2D{1536} [wh] 2D{3072} 4U \\
job(12)/res\_start\_r                & 16L 2H 4L \\
wr\_en\_o                            & 6L 6H 2L 6H 2L \\
wr\_addr\_o                          & 6U [gr] 2D{1024} 2D{2048} 2D{4608} 2U 2D{5120} 2D{1536} [wh] 2D{3072} 2U \\
wr\_data\_o                          & 6U [gr] 2D{15} 2D{13} 2D{10} 2U 2D{10} 2D{13} [wh] 2D{18} 2U \\
\extracode
\begin{pgfonlayer}{background}
   \vertlines[densely dotted, black!50]{0,2,...,22}
\end{pgfonlayer}
\foreach \x/\l in {0/r-1, 2/r, 4/r+1, 6/r+2, 8/r+3, 10/c, 12/c+1, 14/c+2, 16/c+3, 18/c+4, 20/c+5}
   \node[above, font=\small] at (\x+1, 0.9) {\l};
\end{tikztimingtable}
\end{document}
